\section{Limitations}

This study is subject to several limitations that should be considered when interpreting its findings.

First, the systematic review relies exclusively on secondary data sources—published literature and publicly available documentation—which may introduce reporting bias. Studies reporting successful implementations may be more likely to be published and to receive institutional documentation, potentially overstating the prominence of facilitating factors relative to barriers.

Second, the multi-case analysis is limited to four publicly documented national implementations. While these cases were selected to maximize contextual variation, they represent politically prominent programs with extensive public reporting. Less visible implementations—particularly those that experienced failure or were quietly discontinued—are underrepresented, which may introduce survivor bias into the case-level findings.

Third, the systematic review was restricted to English-language publications. This language bias may limit the representativeness of findings with respect to eHealth adoption experiences in non-English-speaking countries, particularly in low- and middle-income settings where digital health adoption challenges may differ substantively from the contexts represented in this review.

Fourth, the U-T-I-O coding framework involved interpretive judgments and was applied by a single reviewer; absence of independent dual-rating limits inter-rater reliability assessment. Two complementary internal checks were performed. (i) On the 17-source domain-presence coding (U/T/I/O addressed: yes / blank), a 10\% seeded self-audit and a supplementary AI-assisted second-rater check (details in \texttt{analysis/audit-check.md}) each yielded 6/8 (75\%) agreement, with disagreements in opposite directions, suggesting the original threshold sits between stricter and broader interpretations. (ii) On the 16 four-case rubric ratings (H / M / L--M / L per domain), a supplementary AI-assisted second-rater pass (Anthropic Claude Opus 4.7, April 2026; protocol and per-cell reasoning in \texttt{analysis/second-rater-protocol.md}) yielded 6/16 (37.5\%) exact agreement and 15/16 (93.8\%) within-one-band agreement (linear-weighted Cohen's $\kappa$ = 0.27, N = 16); the AI rater was systematically more conservative than the original, and two VA/DoD cells (U and T) were subsequently revised from Medium to Low--Medium based on the second-rater finding that the rubric's explicit decision-rule trigger language was more closely matched by the AI's reading. Independent dual-rating with a credentialed second rater remains an outstanding methodological gap for both checks.

Fifth, the breadth of the U-T-I-O Model—spanning user, technical, institutional, and organizational domains—means that individual domain constructs are operationalized at a relatively high level of abstraction. Future research should develop and validate sub-domain measurement instruments to enable more precise quantitative assessment of adoption readiness.

Despite these limitations, the study provides a robust multi-source synthesis of adoption determinants across diverse health system contexts and offers a theoretically grounded, practically oriented framework for guiding eHealth adoption strategy.
